\documentclass[11pt,oneside,openany]{book}
\usepackage[paperwidth=6.5in,paperheight=9.5in,inner=.72in,outer=.72in,top=.75in,bottom=.78in,headheight=14pt]{geometry}
\usepackage{fontspec}
\setmainfont{lmroman10-regular.otf}[BoldFont=lmroman10-bold.otf,ItalicFont=lmroman10-italic.otf,BoldItalicFont=lmroman10-bolditalic.otf,SmallCapsFont=lmromancaps10-regular.otf]
\usepackage[main=latin]{babel}
\usepackage{amsmath,amssymb,graphicx,microtype,fancyhdr,needspace}
\usepackage[unicode,hidelinks]{hyperref}
\hypersetup{pdftitle={Carl Friedrich Gauss: Werke, Band VII. Titelblätter, Praefatio, Contenta und Liber I, 1–114.},pdfauthor={Carl Friedrich Gauss}}
\setlength{\parindent}{1.2em}\setlength{\parskip}{0pt plus .3pt}
\setlength{\emergencystretch}{1.4em}
\pagestyle{fancy}\fancyhf{}
\fancyhead[L]{\small THEORIA MOTUS CORPORUM COELESTIUM}
\fancyfoot[C]{\thepage}\renewcommand{\headrulewidth}{.25pt}
\newcommand{\SourceLeaf}[1]{\hypertarget{native-#1}{}\label{native-#1}}
\renewcommand{\thefootnote}{*)}
\widowpenalty=10000\clubpenalty=10000\displaywidowpenalty=10000
\raggedbottom
\newcommand{\BareRoot}[2][]{\mathord{{}^{#1}\!\surd}#2}

\makeatletter
\newcommand{\HistoricalNoteEight}[1]{\begingroup
\renewcommand{\thefootnote}{[*)]}
\long\def\@makefntext##1{\noindent [*)\ ##1}
\footnote{#1}\endgroup}
\makeatother

\makeatletter
\newcommand{\HistoricalDisplayNoteMark}{\begingroup\renewcommand{\thefootnote}{[*)]}\footnotemark\endgroup}\newcommand{\HistoricalDisplayNoteText}[1]{\begingroup\renewcommand{\thefootnote}{[*)]}\long\def\@makefntext##1{\noindent [*)\ ##1}\footnotetext{#1}\endgroup\ignorespaces}
\makeatother
\newfontfamily\AstroSymbols{Segoe UI Symbol}
\newcommand{\AscendingNode}{\text{\mdseries\AstroSymbols ☊}}

\makeatletter
\newcommand{\HistoricalNoteDouble}[1]{\begingroup
\renewcommand{\thefootnote}{[**)]}
\long\def\@makefntext##1{\noindent [**)\ ##1}
\footnote{#1}\endgroup}
\makeatother

\newcommand{\HistoricalNoteDoubleUnbracketed}[1]{\begingroup\renewcommand{\thefootnote}{**)}\footnote{#1}\endgroup}

\makeatletter
\newcommand{\HistoricalUnmarkedNote}[1]{\begingroup
\renewcommand{\thefootnote}{}
\long\def\@makefntext##1{\noindent ##1}
\footnote{#1}\endgroup\ignorespaces}
\makeatother

\newcommand{\GaussCurrentLastArticle}{114}
